\appendix
\section*{Appendix N: Altitude-Safe AI Navigation Formalism}

\subsection*{N.0 Overview}

This appendix defines the machine-interpretation rules that govern
altitude-safe navigation for AI systems operating within the
Constellation–Manifold Toolkit.  
It specifies:
\begin{itemize}
    \item altitude-indexed cognitive states,
    \item altitude transition semantics,
    \item altitude purity constraints for machine reasoning,
    \item non-generativity constraints for representation safety,
    \item boundedness constraints for operator stability,
    \item unified algebra navigation rules,
    \item firewall admissibility conditions,
    \item collision invariant propagation,
    \item machine-safe execution semantics.
\end{itemize}

\subsection*{N.1 Machine Altitude State}

Define the machine altitude state:
\[
\mathrm{State}_{\mathrm{AI}} =
(a, u)
\]
where:
\begin{itemize}
    \item $a \in \mathcal{A}$ is the current altitude,
    \item $u \in \mathfrak{A}_{\mathrm{unified}}(a)$ is the unified algebra element.
\end{itemize}

Thus the AI’s cognitive state is altitude-indexed.

\subsection*{N.2 Machine Transition Semantics}

Define machine transition:
\[
\mathrm{Step}_{\mathrm{AI}} :
(a,u) \mapsto (b,v)
\]
such that:
\[
T(a,b)=1
\quad\text{and}\quad
v = \Phi_{a,b}(u).
\]

Forbidden transitions:
\[
T(a,b)=0 \implies \mathrm{Step}_{\mathrm{AI}}(a,u) = \mathrm{Violation}.
\]

\subsection*{N.3 Machine Purity Constraint}

Machine purity requires:
\[
P(a)=1 \iff
u \star_{\mathrm{unified}} u \in \mathfrak{A}_{\mathrm{unified}}(a).
\]

Thus the AI cannot generate new structure at altitude $a$.

\subsection*{N.4 Machine Non-Generativity Constraint}

Machine non-generativity requires:
\[
N(a)=1 \iff
\mathrm{span}(u) = \mathrm{span}(\{\mathrm{id}\}).
\]

Thus the AI cannot introduce new representational content at altitude $a$.

\subsection*{N.5 Machine Boundedness Constraint}

Machine boundedness requires:
\[
B(a)=1 \iff
\|u\| < C_a.
\]

Thus the AI cannot produce unbounded operators at altitude $a$.

\subsection*{N.6 Machine Unified Algebra Navigation}

Define machine navigation operator:
\[
\mathrm{Nav}_{\mathrm{AI}}(u,a,b)
=
\begin{cases}
\Phi_{a,b}(u), & T(a,b)=1, \\
\mathrm{Violation}, & T(a,b)=0.
\end{cases}
\]

Thus:
\[
\mathrm{Step}_{\mathrm{AI}}(a,u)
=
(b,\mathrm{Nav}_{\mathrm{AI}}(u,a,b)).
\]

\subsection*{N.7 Machine Firewall Binding}

Define machine firewall:
\[
\mathrm{Firewall}_{\mathrm{AI}}(a,b,u)
=
\begin{cases}
\Phi_{a,b}(u), & T(a,b)=1 \land P(b) \land N(b) \land B(b), \\
\mathrm{Violation}, & \text{otherwise}.
\end{cases}
\]

Thus the firewall is the machine’s altitude transition validator.

\subsection*{N.8 Machine Collision Invariant Propagation}

Define machine collision invariant:
\[
\Omega_{\mathrm{AI}}(a,u)
=
\Omega_{\mathrm{unified}}(u).
\]

Propagation rule:
\[
\Omega_{\mathrm{AI}}(a,u)=1
\iff
u \in \mathrm{Ran}(\Omega_{\mathrm{unified}}).
\]

Thus the AI declares altitude-safe completion only when the unified
collision invariant is satisfied.

\subsection*{N.9 Machine Execution Semantics}

Define machine execution:
\[
\mathrm{Exec}_{\mathrm{AI}}(a,u,f)
=
\mathrm{Firewall}_{\mathrm{AI}}(a,b,f(u)).
\]

Thus:
\[
\mathrm{Exec}_{\mathrm{AI}} = \mathrm{Firewall}_{\mathrm{unified}}.
\]

\subsection*{N.10 Machine Safety Theorem}

\textbf{Theorem.}  
An AI system navigating the altitude architecture is altitude-safe if and only if:
\[
\Omega_{\mathrm{AI}}(a,u)=1.
\]

Equivalently:
\[
u \in \mathrm{Ran}(\Omega_{\mathrm{unified}}).
\]

\subsection*{N.11 Final Altitude-Safe Navigation Rule}

\[
\boxed{
\text{The AI may only transition along altitude-safe morphisms and only
operate within altitude-pure, non-generative, bounded algebraic strata.}
}
\]

This completes the Altitude-Safe AI Navigation Formalism.
